\subsection{无对称性破缺时的解}\label{solutions}

现在我们要用刚刚描述的想法来确定超对称方程 (\ref{commuting})、(\ref{momentmap}) 的某些解模空间。我们将对 $G=SO(3)$ 做这件事，于是用其表示来标记奇异磁单极子的对偶群是 $G^\vee=SU(2)$。

我们首先在库仑分支的原点处工作；这意味着我们考虑标量场 $\vec\phi$ 在 $y\to\infty$ 时为零的解。特别地，$\varphi$ 在无穷远处必须为零。由于方程 $\D_y\varphi=0$ 意味着 $\varphi$ 的共轭类与 $y$ 无关，$\varphi$ 只有在处处幂零时才能在无穷远处为零。这意味着通过一个复规范变换，我们可以使 $\varphi$ 成为上三角的：
\begin{equation}\label{uptri}\varphi=\begin{pmatrix} 0 & p(z) \cr 0 & 0 \end{pmatrix},\end{equation}
其中 $p(z)$ 是某个多项式。然而，我们不能同时把 $s$ 也化成标准形式。因此我们简单地取
\begin{equation}\label{nuptrix}s=\begin{pmatrix} P(z)\cr Q(z)\end{pmatrix},\end{equation}
其中 $P,$ $Q$ 是多项式。在不改变 $\varphi$ 的前提下，我们可以作一个上三角规范变换，把 $P$ 移动 $Q$ 的多项式倍：
\begin{equation}\label{uptrix} P(z)\to P(z)+U(z) Q(z).\end{equation}
除了把 $s$ 乘以一个复常数之外，这是唯一的自由度。

现在我们来问：如何选取 $p, $ $Q$, 和 $P$ 以描述边界上具有磁荷 $k_a$、位于点 $z_a$ 的奇异磁单极子的位形？由于 $s\wedge \varphi s = pQ^2$，我们需要的条件是
\begin{equation}\label{factors}
p(z) Q(z)^2 = \prod_a(z - z_a)^{k_a} := K(z).
\end{equation}
若所有奇异磁单极子都带最小磁荷 $k_a=1$，我们只能以 $p=K$ 与 $Q=\pm 1$ 满足此式。于是可以通过变换 (\ref{uptrix}) 置 $P(z)=0$，而改变 $Q$ 的符号只是把 $s$ 乘以一个无关紧要的常数。所以在库仑分支原点处，对于给定位置与最小磁荷的奇异磁单极子，解是唯一的。
% [Tutor Note: 物理图像] 关键一步是把边界数据压缩为 K(z)=∏(z−z_a)^{k_a}：它只记录奇异磁单极子的位置与电荷，而"把 K 的因子如何分配给 p 与 Q"才编码真正的自由度。原点处 k_a=1 时 K 只能整体给 p、Q 为常数，故解唯一——这正是下文所谓"障碍"的来源：解唯一就意味着只有一维 Hilbert 空间，无法容纳辫群的表示。

这一事实实际上对第~\ref{methods}~节所描述的方案构成一个障碍（见图 2）。最小磁荷奇异磁单极子的情形被认为对偶于 Jones 多项式，而后者是在 $G^\vee=SU(2)$ 且所有纽结都标以 $SU(2)$ 的二维表示时的 Chern-Simons 规范理论中所计算的不变量。若只有一个解，这意味着与此问题相联系的物理 Hilbert 空间是一维的。于是辫群的 Jones 表示将是秩 1 的。这当然不符合事实。我们一定是遇到了图 4 所暗示的某种问题。在把对称性破缺纳入考虑（第~\ref{symbr}~节）以及在形变到 $\t\not=1$（第~\ref{analog}~节）之后，我们会更清楚地看到发生了什么。这最终将使我们能够绕过刚刚描述的障碍。

若某些奇异磁单极子带非最小磁荷，情况会怎样？此时我们会遇到解的模空间。一般地，我们通过
\begin{equation}\label{formsum} p(z)=
\prod_a(z-z_a)^{m_a},~~ Q=\prod_a(z-z_a)^{r_a},\end{equation}
来解 (\ref{factors})，其中
\begin{equation}\label{tonf}m_a+2r_a=k_a.\end{equation}
对 (\ref{formsum}) 中的 $Q$，我们不能通过变换 (\ref{uptrix}) 把 $P$ 置为 1；相反，$P$ 的取值及其前 $r_a-1$ 阶导数在每个点 $z=z_a$ 处都是不变量。因此每个 $z_a$ 联系着 $r_a$ 个模，其中 $r_a$ 的可能取值为 $\{0,1,\dots, [k_a/2]\}$。我们将从第~\ref{symbr}~节开始、并在第~\ref{supermon}~节中最详细地论证：这些模表示某些光滑 BPS 磁单极子在 $y$ 方向上的位置，以及一些共轭角度。
% [Tutor Note: 含糊点拨] "这些模表示 y 方向上的位置"目前只是提议（作者下文也说的是 propose），其确切含义要求在 y 很大、Bogomolny 近似成立时取极限；模的数目 r_a 可由 P 及其前 r_a−1 阶导数在每个 z_a 处的独立取值数出，但要把它们真正解释为位置，需要引入超势并在后文（标签 supermon 的那一节）中定量化。阅读时不要把它当作已证明的定理。

在实施第~\ref{methods}~节所描述的方案时，若某些纽结标以整数 $k_a>1$——换言之，标以 $G^\vee=SU(2)$ 的 $k_a+1>2$ 维表示——人们将不得不把四维解的演化当作三维解模空间中的一条路径来处理。（这可以通过量子化模空间以获得适当的物理态空间来完成，演化就在其中发生。）我们将不走这条路线，而是把刚刚分析的方程扰动到更一般的情形：在无穷远处带上对称性破缺，或取 $\t\not=1$。这最终将把我们所有的模空间约化为有限个点的集合。

\subsection{“实”对称性破缺}\label{symbr}

我们更清楚地理解解——以及更有效地约化到三维——的第一种方法，是沿库仑分支移动。基本想法已在第~\ref{review}~节末尾解释过。我们为切向标量场 $\vec \phi$ 固定某个常数、非零的期待值 $\vec \ca$，并且只考虑超对称方程 (\ref{bpseqns}) 的解中满足 $\lim_{y\to \infty}\vec\phi=g \vec \ca g^{-1}$（对某个 $g\in G$）的那些。

从我们现在的观点看，我们想在对称性破缺的背景下分析约化后的三维方程。约化分离出 $x^1$ 方向，因此我们写  $\vec\phi\cdot \vec \d x=\phi_1\,\d x^1+\varphi\, \d z +\bar \varphi\,\d \bar z$。$\phi_1$ 的期待值或 $\varphi$ 的期待值所产生的影响将颇为不同。我们先考虑只开启 $\phi_1$ 的情形。我们将看到，这会使解在物理上更加透明，而不会增加或减少解。

对 $G=SO(3)$，我们可以选取一个规范，使 $\phi_1$ 在大 $y$ 处为常数。这意味着在大 $y$ 处，$\A_y=A_y-i\phi_1$ 在无穷远处趋于一个常数矩阵，例如使得
\begin{equation}\label{gauga}
\A_y = \begin{pmatrix} \ca_1 & 0 \cr 0 & -\ca_1\end{pmatrix}.
\end{equation}
其中 $\ca_1>0$。这一期待值把 $SO(3)$ 破缺为 $U(1)$。我们选取未破缺 $U(1)$ 的一个生成元，其归一化使得伴随表示值场的非对角分量（对 $G=SO(3)$ 它们是最小电荷的场）具有电荷 $\pm 1$：
\begin{equation}\mathcal Q=\begin{pmatrix}-1/2 & 0 \cr 0 & 1/2\end{pmatrix}\end{equation}
符号的选取将是方便的。
% [Tutor Note: 约定差异] 生成元归一化为 Q=diag(−1/2,1/2)，才使非对角（即最小电荷）分量的电荷为 ±1；若改用 Q=diag(−1,1)，则最小电荷变为 ±2 而光滑磁单极子的磁荷仍为 2，所有"每加一个零点磁荷增 2"的表述都要随之调整。文中"符号的选取将是方便的"正是指这一约定，另注意 ca_1>0 的取法只是为了固定符号。

方程 $\D_y\varphi=0$ 意味着 $\varphi$ 的下三角分量必须为零，否则它会在无穷远处指数式增长。对角分量也必须为零，否则 $\varphi$ 在无穷远处会有非零极限。因此对称性破缺提供了一个自然的标架，使 $\varphi$ 在 $y\to\infty$ 时严格上三角，正如边界条件在 $y\to 0$ 处所做的那样。但 $\varphi$ 在大 $y$ 处指数式衰减。

在低能有效 $U(1)$ 理论中，电荷为 1 的场是我们将称之为 $\M$ 的线丛的一个截面。把 $\M$ 的第一 Chern 类在由 $y=y_0$（$y_0$ 为某个大常数）给出的平面上积分，得到一个整数，我们称之为磁荷 $\m$。$\varphi$ 的上三角矩阵元
\begin{equation}\label{tomigo}\varphi\sim \begin{pmatrix} 0 & \tilde p \cr 0 & 0 \end{pmatrix}
\end{equation}
是 $\M^{-1}$ 的一个截面，因此磁荷（即 $\M^{-1}$ 的第一 Chern 类的负值）等于 $\tilde p$ 的零点数的负值。这与 $p$（定义于 (\ref{uptri})）的零点数的负值相同，而后者通过一个复规范变换等价于 $\tilde p$。于是在 (\ref{tonf}) 的记号下，磁荷为
\begin{equation}\label{magch} \m = -\sum_a m_a=\sum_a(-k_a+2r_a).\end{equation}
这意味着，对于 $y=0$ 处给定的奇异磁单极子或 't Hooft 算符位形，$\m$ 的最小可能值是 $-\sum_ak_a$。每当我们给 $Q$ 增加一个零点，$\m$ 就增加 2。
% [Tutor Note: 易错点] 磁荷的符号约定容易出错：文中取 m 为线丛 M 的第一 Chern 类的负值，而 φ 的上三角元 p̃ 是 M^{-1} 的截面，于是"磁荷 = −p̃ 的零点数"。若把 M 与 M^{-1} 弄反，式 (\ref{magch}) 的整体符号以及"最小磁荷为 −∑k_a"的结论都会颠倒。另注意光滑磁单极子的磁荷为 2 而非 1，这正是本段所有"每加一个 Q 的零点 m 增 2"的来源。

我们提议这一事实可以用光滑 BPS 磁单极子来解释。若简单地令 $\varphi=0$，我们一直研究的推广 Bogomolny 方程就约化为通常的 Bogomolny 方程。远离边界处，由于 $\varphi$ 如此之小，Bogomolny 方程是极好的近似。$\R^3$ 上的 Bogomolny 方程允许磁荷为 $2$ 的光滑磁单极子解。（我们测量磁荷所用的单位，使得 Dirac 量子化所允许的最小磁荷为 $\pm 1$。这也正是最小磁荷奇异磁单极子或 't Hooft 算符的磁荷，但光滑磁单极子具有偶数磁荷。）我们的提议是：若 $Q$ 在 $z=z_a$ 处有 $r_a$ 阶零点，则有 $r_a$ 个光滑磁单极子位于 $z=z_a$。这一陈述只有在磁单极子位于很大的 $y$ 处（那里推广 Bogomolny 方程约化为通常方程并以光滑磁单极子为解）时才有确切含义。然而我们发现，若 $Q$ 在 $z=z_a$ 处有 $r_a$ 阶零点，则在复 $z$ 平面中与点 $z=z_a$ 相联系的复模恰有 $r_a$ 个。我们提议这些模就是位于 $z=z_a$ 的 $r_a$ 个光滑磁单极子在 $y$ 方向上的位置，外加共轭角度。同样，这一陈述的确切含义在所涉 $y$ 位置很大时成立。我们将在第~\ref{supermon}~节中定量地建立这一解释，但现在只考虑定性的论证。

这一提议的第一个动机来自 $\m$：我们把公式 (\ref{magch}) 解释为边界处奇异磁单极子的磁荷之和（$-\sum_ak_a$）加上每个光滑磁单极子贡献的 2（总共贡献 $2\sum_ar_a$）。奇异磁单极子的磁荷为负这一点值得注意。在对称性破缺之前，奇异磁单极子的磁荷没有有意义的符号——它相对于 $G^\vee=SU(2)$ 的一个表示来说是"对偶"的——但在对称性破缺之后，这个符号有了意义，而在我们的归一化下它是负的。


为获得更多知识，我们回顾关于 Bogomolny 方程的一些事实。对 $\R^3$ 上的 Bogomolny 方程，全纯数据就是对易算符 $\D_1$ 与 $\D_3$，其定义与式 (\ref{neqform}) 完全相同。Bogomolny 方程的定域磁单极子解通过 $\D_3$ 平行移动中的“束缚态”体现在全纯数据中，即通过
\begin{equation}\label{zonurf}\frac{\D s}{\D y}=0.\end{equation}
的可归一化解体现出来。这样的解只出现在 $z$ 平面中的特定位置，这些位置被解释为磁单极子位置的 $z$ 值，因为一般地，若我们取 $s$ 在 $y\to -\infty$ 时指数衰减，则它在 $y\to +\infty$ 时将指数增长。

在我们的设定中，我们局限于 $y\geq 0$，而类似的归一一化解恰好存在于如下情形：在 $y\to 0$ 时消失的“小”截面 $s$ 在大 $y$ 处也指数衰减。若 $\varphi$ 同样指数衰减，则 $s\wedge \varphi s$ 在 $y\to\infty$ 时为零，从而对所有 $y$ 为零。所以两端都消失的截面 $s$ 只能出现在 $K(z)$ 的零点处。的确，若在无穷远处附近
\begin{equation}\label{sinfty}s=\begin{pmatrix} \tilde P\cr \tilde Q\end{pmatrix}\end{equation}
则给定 (\ref{gauga}) 的形式，$s$ 在 $y\to \infty$ 时消失等价于 $\tilde Q=0$。$\tilde Q$ 的零点与 $Q$ 的零点相同。
（的确，$s\wedge \varphi s$ 与 $y$ 无关；其零点在 $y$ 小时是 $Q$ 的零点，在 $y$ 大时是 $\tilde Q$ 的零点。）
因此，与通常 Bogomolny 方程的关系确实提示：在每个使 $Q$ 为零的点 $z=z_a$ 处有 $r_a$ 个光滑磁单极子。光滑磁单极子在 $y$ 方向上的位置应当编码在 $P$ 及其前 $r_a-1$ 阶导数的取值中。我们将在第~\ref{supermon}~节中使这一图景更精确；现在只需注意，可归一化的小截面 $s(z_a)$ 的行为是 $P(z_a) e^{ -m_1 y}$，这提示增大 $|P(z_a)|$ 会把光滑磁单极子推向大 $y$。
% [Tutor Note: 原文笔误] 原文此句为 "can arise only only at zeroes of K(z)"，重复了 only 一词，属英文笔误；译文按单次"只能"处理。此外"两端都消失的截面只在 K(z) 的零点处出现"这一断言的严格论证，依赖于下一句关于 s∧φs 与 Q（或 Q̃）零点一致的说明。

极限 $P(z_a) \to 0$ 有些奇异，因为 $s$ 本应处处非零。物理图景提示该极限将对应于磁单极子起泡（monopole bubbling）的情形：光滑磁单极子被推到边界，并屏蔽奇异磁单极子的磁荷。磁单极子起泡是一种现象（最初由 \cite{Kronheimer} 描述，并在 \cite{KW} 第 10.2 节中被重新发现），其中 't Hooft 算符吸收一个光滑 BPS 磁单极子，从而降低其磁荷的大小。
% [Tutor Note: 物理图像] 磁单极子起泡：当 P(z_a)→0 时小截面 s 的归一化失效，光滑磁单极子被推到 y=∞ 并被 't Hooft 算符吸收，后者磁荷由 −1 降为 +1（即减少 2）。这不是数学上的退化，而是屏蔽/隧穿的物理体现；这一图像负责解释下文"把 K 的零点任意分配给 P 与 Q"时出现的共零点情形。

\subsection{“复”对称性破缺}\label{simvar}


现在考虑 $\varphi$ 中的对称性破缺情形。设对 $G=SO(3)$，$\varphi$ 在 $y\to\infty$ 时的本征值为 $\pm \ca$。于是，由于 $\varphi$ 与 $\D_y$ 对易，其本征值处处为 $\pm \ca$。通过一个复规范变换，我们可以化到 $\varphi$ 为常数对角矩阵的情形：
\begin{equation}\label{constdiag}\varphi=\begin{pmatrix} \ca & 0 \cr 0 & -\ca \end{pmatrix}.\end{equation}
照常把“小”截面写成 $s=\begin{pmatrix}P \cr Q\end{pmatrix}$，
我们得到 $s\wedge \varphi s = 2\ca  PQ$。在把 $\varphi$ 化为形式 (\ref{constdiag}) 之后，
我们仍可作 $\mathrm{diag}(\lambda,\lambda^{-1})$ 的规范变换，它把
\begin{equation}\label{omog}P\to \lambda P,~~Q\to\lambda^{-1} Q.\end{equation}
在 $z=z_a$（及 $y=0$）处存在磁荷为 $k_a$ 的奇异磁单极子时，
我们所需满足的条件是 $s\wedge\varphi s=\prod_{a=1}^\ss (z-z_a)^{k_a}:=K(z)$。
这成为
\begin{equation}\label{counting}
2\ca P Q  = K(z).
\end{equation}
这些方程的解与 $K(z)$ 的分解相联系，并且（模去变换 (\ref{omog})）没有模。

在所有 $k_a=1$（对应最小磁荷）的 $\ss$ 个边界 't Hooft 算符的情形下，解的数目恰为 $2^\ss$。这些解简单地对应于把 $K(z)$ 的因子分配给 $P$ 与 $Q$ 的各种方式。数目 $2^\ss$ 有一个自然的解释。在库仑分支上，我们可能期望最小磁荷的 't Hooft 算符有两种可能的状态：其磁荷与对称性破缺方向同向或反向。这两个状态对应于 $P$ 中的零点或 $Q$ 中的零点。在对偶的、带 Wilson 算符的 Chern-Simons 理论描述中，最小磁荷的 't Hooft 算符对应于 $SU(2)$ 二维表示中的 Wilson 算符；这样的算符同样表示两个量子态，其沿对称性破缺轴的电荷为正或为负。

为确认这 $2^\ss$ 个解对应于每个 't Hooft 算符磁荷的两种可能选择，让我们计算这些解的磁荷。设 $P$ 的次数为 $\ss_1$，$Q$ 的次数为 $\ss_2$，其中 $\ss_1+\ss_2=\ss$。比值 $P/Q$ 不依赖于小截面 $s$ 的归一化。在低能阿贝尔规范理论中，该比值具有电电荷 $-1$；它是线丛 $\M^{-1}$ 的一个截面。另一方面，具体地说，$P/Q$ 在 $z$ 平面上有 $\ss_1$ 个零点与 $\ss_2$ 个极点。（把 $P/Q$ 理解为 $\M^{-1}$ 的截面时，它在 $z=\infty$ 处既无极点也无零点。事实上，通过一个复规范变换，可以把解在 $z=\infty$ 附近化成原始 Nahm 极解 (\ref{nahmpole}) 的形式，特别地它与 $z$ 无关。这使 $\M$ 在 $z=\infty$ 附近平凡化，并使 $P/Q$ 与 $z$ 无关。）于是线丛 $\M$ 的次数为 $\ss_2-\ss_1$，从而
\begin{equation}\label{melf}\m = \ss_2-\ss_1.\end{equation}
由此可见，$z=z_a$ 处的 't Hooft 算符对 $\m$ 贡献 $-1$ 或 $+1$，取决于我们把 $z-z_a$ 的因子放进 $P$ 还是 $Q$。

为理解为什么 $\ca\not=0$ 时的解比 $\ca=0$ 时多，考虑 $\varphi$ 的一个略更一般的拟设：
\begin{equation}\label{pony}\varphi=\begin{pmatrix} \ca & p(z)\cr 0 & -\ca\end{pmatrix}.
\end{equation}
对 $\ca\not=0$，多项式 $p(z)$ 可以通过上三角规范变换去掉，但现在 $\ca=0$ 时存在光滑极限。
 条件 $s\wedge \varphi s=K$ 成为
\begin{equation}\label{turbo}  2\ca PQ+pQ^2 = \prod_{a=1}^\ss(z-z_a).\end{equation}
对 $\ca\not=0$，$p$ 无关紧要，因为它可以通过 $P\to P-pQ/2\ca$ 消去。
对 $\ca\not=0$，方程 (\ref{turbo}) 有 $2^\ss$ 个解（模去保持 $\varphi$ 形式的复规范变换），对应于把 $K(z)$ 分解为 $(2\ca P+pQ)Q$。但这些解都有 $P\sim 1/\ca$，只有我们已在第~\ref{solutions}~节中找到的那个 $Q=1$（且 $P=0$）的解例外。

我们对此作如下解释。
在边界点 $z=z_a$ 处存在最小磁荷 't Hooft 算符时，总有一个局域于边界附近的磁荷 $-1$。
局域于边界附近的磁荷对所有解都相同，因为解在边界附近的形式总是由标准模型解（\cite{fiveknots} 第 3.6 节所描述的那个）给出，
与其余一切无关。在 (\ref{counting}) 的一个满足 $Q(z_a)=0$ 且（因而）
$P(z_a)\not=0$ 的解中，另外还有一个磁荷为 $2$ 的光滑 BPS 磁单极子，位于
$z=z_a$ 及某个依赖于 $P(z_a)$ 的 $y$ 值处。所以与 $z=z_a$ 相联系的总磁荷为 $-1+2=1$。当 $\ca\to 0$ 时，$P(z_a)\to\infty$，
光滑磁单极子消失到 $y=\infty$。这样，所有 $2^\ss$ 个解在 $\ca\to 0$ 时都变得等价，尽管它们在 $\ca\not=0$ 时互不相同。一如既往，这种用光滑 BPS 磁单极子给出的描述，只有在磁单极子位于很大的 $y$ 值处、从而使 Bogomolny 方程在其位置附近是好的近似时才精确。正如我们将在第~\ref{supermon}~节中最充分地说明的，这恰好发生在 $\ca$ 非常小（与 't Hooft 算符之间的倒数距离 $1/(z_a-z_b)$ 相比）的时候，所以
特别地，BPS 磁单极子描述在 $\ca\to 0$ 时变得精确。

在相反的情形，即 $\ca$ 与倒数距离相比很大时，对称性破缺很强，边界上不同的奇异磁单极子相距如此之远，以致它们不会显著地相互影响。在此情形下，本质的陈述只是：一个最小磁荷的 't Hooft 算符，在 $\ca\not=0$ 的对称性破缺存在时，
有两个可能的状态，其中它在有效低能理论中表现为电荷 1 或 $-1$ 的 't Hooft 算符。据此，
由 $\ss$ 个相距很远的最小磁荷 't Hooft 算符构成的系统自然地有 $2^\ss$ 个可能状态。


\subsubsection{推论}\label{physhilb}

对 $\ca\not=0$ 时解的计数绕过了我们在第~\ref{solutions}~节中遇到的困难。有了 $2^\ss$ 个经典解，在存在 $\ss$ 个最小磁荷 't Hooft 算符时的物理 Hilbert 空间将有维数 $2^\ss$。这 $2^\ss$ 个态有可能对偶于标以 $SU(2)$ 二维表示的 $\ss$ 个 Wilson 算符的态。而与 Jones 多项式相联系的、$\ss$ 股辫群的表示，当然可以在一个 $2^\ss$ 维的向量空间中实现。

据此，有理由期望：若 $\ca$ 是一般的，第~\ref{methods}~节图~\ref{stretching} 所勾勒的拉伸策略可以奏效；但对 $\ca=0$，试图通过沿一个方向拉伸来简化经典解，会导向图~\ref{failure} 所暗示的那类行为。

更一般地，在边界上存在磁荷为 $k_a$ 的 't Hooft 算符时，我们希望找到磁荷为 $-k_a,-k_a+2,-k_a+4,\dots,k_a$ 的 $k_a+1$ 个态。的确，$SO(3)$ 规范理论中的这样一个 't Hooft 算符对偶于 $SU(2)$ 规范理论中与自旋 $k_a/2$ 的表示相联系的 Wilson 算符；该表示的维数为 $k_a+1$，其权正如所指出的那样。如果在解方程 (\ref{counting}) 时允许 $K(z)$ 的任意分解、把零点以任意方式分配给 $P$ 与 $Q$，我们确实能得到正确数目的、具有正确磁荷的解。不幸的是，对于 $P$ 与 $Q$ 都在 $z=z_a$ 处为零的那些分解，我们没有简单的解释。毕竟 $s=\begin{pmatrix}P \cr Q\end{pmatrix}$ 本应处处非零。看起来有可能：$P$ 与 $Q$ 有公共零点的分解应当用磁单极子起泡来解释，即解释为在存在降低后磁荷为 $k_a-2$、$k_a-4$ 等的 't Hooft 算符时方程的解。详细发展这一想法超出了本文的范围。

回到所有 $k_a$ 都等于 1 的较简单情形：由于物理态空间的维数 $2^\ss$ 很有希望，下一步可以是尝试计算辫群的作用并提取 Jones 多项式。在尝试这样做之前，我们将描述问题的进一步形变，它将被证明是有启发性的。除其他之外，通过这一形变，无论 't Hooft 算符的磁荷如何，我们都会得到良好的行为。
